\section{Household Survey in Argentina}\label{app:EPH}\setcounter{equation}{0}


The Encuesta Permanente de Hogares (EPH) is Argentina’s main household survey for tracking labour market conditions and broader socio-economic indicators in urban areas. Conducted by the Instituto Nacional de Estadística y Censos (INDEC), its primary objective is to provide timely information on employment, unemployment, underemployment, hours worked, and income, along with complementary data on education, housing, migration, and other household characteristics. The EPH is a continuous survey with quarterly frequency, which makes it the cornerstone of official labour market statistics in Argentina and the basis for the calculation of indicators such as the unemployment rate and urban poverty measures.

A distinctive feature of the EPH is its rotation panel design. Each selected household is interviewed four times: in two consecutive quarters of a given year, and then again in the same two quarters of the following year. For example, a household first surveyed in quarter 1 will also be surveyed in quarter 2, and then again in quarters 5 and 6 (the same quarters of the subsequent year). This structure generates a short panel in which roughly half of the sample can be tracked over time, enabling researchers to study both short-term (quarter-to-quarter) and longer-term (year-over-year) transitions in labour market status. Such a design balances the need for fresh samples to ensure representativeness with the ability to follow the same households over time.

The survey is restricted to urban agglomerates—currently covering 31 metropolitan areas across the country—which together represent a large share of Argentina’s population. Within each agglomerate, households are chosen through a stratified, multi-stage sampling design, and data are collected using probabilistic methods to ensure representativeness. The quarterly frequency of the EPH allows for the timely monitoring of economic shocks, structural changes, and seasonal patterns in employment and income, making it particularly useful for both short-run policy responses and longer-run labour market analysis.

We focus on the last two quarters of 2006 because this period coincides with a phase of sustained economic expansion in Argentina, following the recovery from the 2001–2002 crisis but preceding the increase in coverage from the large-scale tax amnesties in this period, and before global turbulence from the Great Recession. By focusing on this pre-crisis, pre-policy-shock moment, we capture a relatively stable baseline against which later changes can be more clearly identified. The survey asks respondents about their contributions to social security. We use this to separate workers into formal and informal, and only use data for those aged 25 to 60.


\subsection{The vector-\texorpdfstring{$X_i$}{Xi} calibration}\label{app:EPH:vectorX}

Section \ref{sec:argentina} identifies the formal households' taste for leisure as one parameter common
to the four quartiles, $X_i = X$, and uses the observed average workweek to fix its level, so that
relative hours across quartiles are a prediction. The alternative is to treat the survey's relative hours
as data at every point of the formal distribution: with both $z_i^{\eta}$ and $z_i^{x}$ imposed,
\eqref{eq:calibration_eta} identifies a full vector $X_i$ and no separate hours target is needed. The
level of average hours is then not identified, and hours can only be read against a reference point.

Table \ref{table:Arg:Calib_vectorX} reports that calibration. Nothing in section \ref{sec:argentina}
depends on which of the two we use, and this is a property of the model rather than a numerical accident:
$X$ enters no aggregate, because the productivity distribution reaches every aggregate through
$\sum_{i>0}\gamma_i\eta_i^{1+\xi}/X_i^{\xi}$ alone, which both variants normalize to one. The calibration
is therefore block recursive, and $\beta$, $\omega$, the capital-output ratio, the equilibrium tax rate,
the savings rate and the informal savings ratio are identical across the two variants at every $\rho$ of
the grid. The reform in table \ref{table:Argentina:Universal}, its dependence on $\epsilon$ and $\theta$
in figure \ref{fig:Argentina_functionOfParams}, and the CRRA results in table
\ref{table:Argentina:funcOfRho} and figure \ref{fig:ARG:EffectOfCRRA} therefore reproduce to every printed
digit, and we do not repeat them. What the variants disagree about is the calibrated productivity and
leisure parameters themselves, and the reading of relative hours: data under the vector-$X_i$ calibration,
and under the common-$X$ one a prediction that reproduces the survey's ranking but overstates the top
quartile's hours by about 8\%.

\input{Tables/ArgentinaCalibration_vectorX}
